%------------------------------------------------------------------------------%

\usepackage{apresentacao}
\usepackage{tabto}

\graphicspath{{./figs/}}

%------------------------------------------------------------------------------%
\newcommand{\mytitle}   {Séries Numéricas com Termos Gerais}
\newcommand{\mysubtitle}{}
\title   {\mytitle}
\subtitle{\mysubtitle}
\author  {\large Luis Alberto D'Afonseca}
\date    {}

%------------------------------------------------------------------------------%
\begin{document}

%------------------------------------------------------------------------------%
\begin{frame}
    \titlepage
\end{frame}

%------------------------------------------------------------------------------%
\section{Séries Alternadas}

%------------------------------------------------------------------------------%
\begin{frame}{Séries Alternadas}
  Os \emph{termos trocam de sinal} alternando entre um positivo e um negativo
  \[ \sum_{n=1}^\infty (-1)^{n+1} a_n \]
  com $a_n >0$ para todo $n$.

  \pause
  \vspace{\baselineskip}
  Exemplo, \emph{Série Harmônica Alternada}
  \[ \sum_{n=1}^\infty \frac{(-1)^{n+1}}{n}  \;=\; 1 - \frac{1}{2} + \frac{1}{3} - \frac{1}{4} + \cdots \]
\end{frame}

%------------------------------------------------------------------------------%
\begin{frame}[t]{Teste da Série Alternada -- Teste de Leibniz}

  \vspace{\baselineskip}
  Seja $\,a_n\,$ tal que, para todo $ n \in \N$,

  \pause
  \vspace{0.5\baselineskip}
  \begin{enumerate}[<+->]
      \item $a_n>0$           \tabto{30mm} $a_n$ são os módulos dos termos da série  \\[2mm]
      \item $a_n\geq a_{n+1}$ \tabto{30mm} os termos são decrescentes em módulo      \\[2mm]
      \item $a_n\to 0$        \tabto{30mm} os termos tendem a zero
  \end{enumerate}

  \only<+->{
  \vspace{1.5\baselineskip}
  Então, a série alternada
  \(\; \sum_{n=1}^\infty (-1)^{n+1}a_n \;\)
  converge.}
\end{frame}

%------------------------------------------------------------------------------%
\begin{frame}[t]{Exemplo}
  Vamos analisar a série
  \[ \sum_{n=1}^{\infty} (-1)^{n+1} \frac{1}{\,2^{n-1}\,} \]
  \pause
  Percebemos que
  \begin{itemize}[<+->]
    \item a série é alternada
    \item o módulo do termo geral é \(\quad a_n = \frac{1}{\,2^{n-1}\,} \)
    \item $a_n\geq a_{n+1}$
    \item $a_n\to 0$
  \end{itemize}
  \only<+->{Portanto a série converge pelo Teste da Série Alternada}
\end{frame}

%------------------------------------------------------------------------------%
\begin{frame}{Exemplo}
  Por outro lado
  \[
    \sum_{n=1}^{\infty} (-1)^{n+1} \frac{1}{\,2^{n-1}\,}
    \pause
    \;=\; \sum_{n=1}^{\infty} \frac{(-1)^2(-1)^{n-1}}{\,2^{n-1}\,}
    \pause
    \;=\; \sum_{n=1}^{\infty} \left( -\frac{1}{\,2\,} \right)^{n-1}
  \]
  \pause
  Série geométrica com $\;a=1\;$ e \(\;r=-\frac{1}{2}\;\) \\[5mm]

  \pause
  \vspace{0.5\baselineskip}
  Como $|r|<1$ ela converge \\[5mm]

  \pause
  sua soma é
  \(\quad
    S \;=\; \frac{a}{1-r}
      \pause
      \;=\; \frac{1}{1+\sfrac{1}{2}}
      \;=\; \frac{2}{3}
  \)

\end{frame}

%------------------------------------------------------------------------------%
\framedgraphic{Usando Planilhas Eletrônicas}{./figs/exemplo_planilha}{}
\framedgraphic{Usando Planilhas Eletrônicas}{./figs/exemplo_grafico}{}

%------------------------------------------------------------------------------%
\section{Estimativa do Erro para Séries Alternadas}

%------------------------------------------------------------------------------%
\begin{frame}{Estimativa do Erro}
  Dada uma série que satisfaz as condições do teste da série alternada
  \pause
  \[ S_n = a_1 - a_2 + a_3 - a_4 + \cdots + (-1)^{n+1}a_n \]
  converge para a soma da série $S$
  \pause
  \begin{itemize}[<+->]
    \item \( \abs{R_n} < a_{n+1} \)
    \item $S$ está entre $S_n$ e $S_{n+1}$
    \item \( R_n = S - S_n \;\) tem o mesmo sinal do primeiro termo não utilizado
  \end{itemize}
\end{frame}

%------------------------------------------------------------------------------%
\section{Convergência Absoluta}

%------------------------------------------------------------------------------%
\begin{frame}{Convergência Absoluta}
  Uma série
  \(\quad \sum_{n=1}^{\infty}\; a_n\) \\[5mm]
  \emph{converge absolutamente}, ou é \emph{absolutamente convergente}, \\[5mm]
  se
  \(\quad\sum_{n=1}^{\infty} \;\abs{a_n} \quad\) converge
\end{frame}

%------------------------------------------------------------------------------%
\begin{frame}{Convergência Condicional}
  Se  \tabto{12mm} \(\sum_{n=1}^{\infty} a_n\)          \tabto{32mm} converge \\[9mm]
  mas \tabto{12mm} \(\sum_{n=1}^{\infty} \;\abs{a_n} \) \tabto{32mm} diverge  \\[9mm]
  a série é \emph{condicionalmente convergente} \\[9mm]
  \pause
  Exemplo: Série Harmônica Alternada
\end{frame}

%------------------------------------------------------------------------------%
\begin{frame}{Teste da Convergência Absoluta}
  Uma série absolutamente convergente é convergente.\\[3mm]

  \pause
  \vspace{\baselineskip}
  Séries com termos não negativos e convergentes
  são absolutamente convergentes
\end{frame}

%------------------------------------------------------------------------------%
\begin{frame}{Teorema do Rearranjo}
Se \(\quad\sum_{n=1}^{\infty} a_n \quad \) converge absolutamente \\[5mm]

\pause
dado qualquer \emph{rearranjo} $(b_k)$ da sequência $(a_n)$ \\[5mm]

\pause
a série $\sum b_k$ converge e
\[
  \sum_{k=1}^{\infty} b_k = \sum_{n=1}^{\infty} a_n
\]
\end{frame}

%------------------------------------------------------------------------------%
%------------------------------------------------------------------------------%
\section{Resumo dos Testes}
%------------------------------------------------------------------------------%
%------------------------------------------------------------------------------%

%------------------------------------------------------------------------------%
\begin{frame}{Testes Básicos}
	\begin{itemize}[<+->]
		\item \emph{Teste da divergência ou do $n$-ésimo termo}\\
          a menos que $a_n\to 0$, a série diverge \\[5mm]
		\item \emph{Séries geométricas} \\[3mm]
          \(\sum_{n=1}^{\infty} ar^{n-1}\) converge se \(\abs{r}<1\); caso contrário, diverge \\[5mm]
		\item \emph{Série $p$} \\[3mm]
          \(\sum_{n=1}^{\infty} \frac{1}{n^p}\)
          converge se $p>1$; caso contrário, diverge\\[5mm]
    \item Série \emph{Telescópica}
	\end{itemize}
\end{frame}

%------------------------------------------------------------------------------%
\begin{frame}{Séries de Termos Não-Negativos}
\begin{itemize}[<+->]
  \item Teste da \emph{integral} \\[5mm]
  \item Teste da \emph{razão} ou o teste da \emph{raiz} \\[5mm]
  \item Teste da \emph{comparação}
\end{itemize}
\end{frame}

%------------------------------------------------------------------------------%
\begin{frame}{Séries Gerais -- Alguns Termos Negativos}
	\begin{itemize}[<+->]
    \item A série é \emph{alternada}? \\[3mm]
          Verifique as condições do teste da série alternada. \\[5mm]
		\item \(\sum_{n=1}^{\infty} \abs{a_n}\) converge? \\[3mm]
		      Se convergir, \(\sum_{n=1}^{\infty} a_n\) também converge
	\end{itemize}
\end{frame}

%------------------------------------------------------------------------------%
\begin{frame}{\mytitle}{\mysubtitle}
  \tableofcontents
\end{frame}

%------------------------------------------------------------------------------%
\end{document}
%------------------------------------------------------------------------------%
